\documentclass[conference]{IEEEtran}

% =========================================================
% PLANTILLA LIGERA PARA OVERLEAF FREE
% Compilador recomendado: pdfLaTeX
% =========================================================

\usepackage[utf8]{inputenc}
\usepackage[T1]{fontenc}
\usepackage[spanish,es-nodecimaldot]{babel}
\usepackage{lmodern}
\usepackage{microtype}

\usepackage{amsmath,amssymb}
\usepackage{graphicx}
\usepackage{booktabs}
\usepackage{multirow}
\usepackage{array}
\usepackage{url}
\usepackage{cite}
\usepackage{xcolor}
\usepackage{listings}
% Se activan cuando el documento incluya algoritmos.
\usepackage{algorithm}
\usepackage{algpseudocode}
\usepackage{hyperref}

\definecolor{CodeBackground}{HTML}{F4F6F7}
\definecolor{CodeKeyword}{HTML}{005F73}
\definecolor{CodeComment}{HTML}{6C757D}
\definecolor{CodeString}{HTML}{BB3E03}

\hypersetup{
    colorlinks=true,
    linkcolor=black,
    citecolor=black,
    urlcolor=blue
}

\lstdefinestyle{paperstyle}{
    backgroundcolor=\color{CodeBackground},
    basicstyle=\ttfamily\footnotesize,
    keywordstyle=\color{CodeKeyword}\bfseries,
    commentstyle=\color{CodeComment}\itshape,
    stringstyle=\color{CodeString},
    numbers=left,
    numberstyle=\tiny,
    numbersep=5pt,
    frame=single,
    breaklines=true,
    showstringspaces=false,
    tabsize=4
}

\lstset{
    style=paperstyle,
    language=Python
}

% =========================================================
% DATOS DEL ARTÍCULO — EDITAR AQUÍ
% =========================================================

\title{
Lectura y analisis de señales ECG
}

\author{
\IEEEauthorblockN{
Kevin Javier Gonzalez Luna,
Iván Felipe Maluche Suarez,
Angélica Pascagaza Vega
}
\IEEEauthorblockA{
\textit{Universidad Nacional de Colombia}\\
\textit{Sede Bogotá}\\
Bogotá D.C., Colombia\\
\{kegonzalezl, imaluche, apascagaza\}@unal.edu.co
}
}
\begin{document}

\maketitle

% =========================================================
% RESUMEN
% =========================================================

\begin{abstract}
Este documento presenta un estudio y posterior implementacion de un lector de señales ECG por medio de las estructuras de datos aprendidas a lo largo del curso, para esto el grupo se basa en diferentes papers y estudios que permitan comprender a fondo la lectura de estas señales y como mostrarlas y leerlas adecuadamente
\end{abstract}

\begin{IEEEkeywords}
Medicina, Electro cardiograma, estructuras de datos,
algoritmos, instrumentos de medicion.
\end{IEEEkeywords}

% =========================================================
% CONTENIDO
% =========================================================

\section{Introducción}

En el campo de la cardiologia debido a la complejidad y gravedad de las enfermades se hace de especial importancia el uso de herramientas capaces de leer y analizar datos clinicos de los pacientes de manera que se pueda advertir de posibles riesgos para su salud, sabiendo como esta tarea puede encontrar solucion en la programacion y las estructuras de datos como grupo vemos de especial interes hacer uso de una implementacion de lo visto en el curso
para clasificar los datos clinicos de un paciente de forma que se puedan leer y analizar de manera eficiente por parte del personal de salud.

Para este fin en nuestro proyecto se estudiara la lectura, analisis y graficacion de señales ECG (obtenidas a base de un electrocardiograma) para ser capaces de entregar un reporte en el que se puedan identificar patrones anomalos.

\subsection{Definicion del proyecto}
La idea central del proyecto consiste en por medio de las implementaciones estudiadas en el curso de estructuras de datos manejar los datos del ECG y detectar anomalias, esto puede ser de gran ayuda pues las estructuras de datos vistas pueden hacer mas eficiente el costo de operacion y por tanto facilitar el trabajo para el personal medico. 
\textbf{[Describir con precisión la situación que se desea resolver, sus
causas, restricciones y consecuencias.]}

\subsection{Objetivos}

\begin{itemize}
    \item \textbf{Ser capaz de leer registros de la base de datos MIT-BIH Arrhythmia Database incluyendo los archivos de registro y las anotaciones del cardiologo.}
    \item \textbf{Usar estructuras de datos capaces de permitir un recibiendo y envio eficiente de datos para el procesamiento y filtrado de informacion}
    \item \textbf{Entregar un resultado entendible que permita analizar los datos}
\end{itemize}
\section{requerimientos basicos de la implementacion}
\label{sec:metodologia}

Esta seccion se enfoca en los requerimientos mas basicos que debe cumplir el codigo.
\\
Para cumplir con su funcion de manera eficiente el programa debe:

\begin{itemize}
 \item Ser capaz de leer cada uno de los archivos del conjunto de datos MIT (.head, .dat, .atr y .xws)
 \item Organizar los datos de manera que puedan ser facilmente leidos y manejados por las estructuras de datos
 \item Ser capaz de graficar los datos de manera que sean entendibles
 \item FIltrar ruido para ver de manera mas clara los patrones de la señal
\end{itemize}

\section{Resultados}
\label{sec:resultados}

Esta sección presenta los resultados obtenidos durante la evaluación de la
solución.

\subsection{Resultados principales}

\textbf{[Presentar los resultados más importantes de forma clara y objetiva.
Evitar interpretar aquí los resultados en exceso.]}

\begin{table}[ht]
\caption{Plantilla para resumir los resultados}
\label{tab:resultados}
\centering
\begin{tabular}{lcc}
\toprule
\textbf{Caso} & \textbf{Métrica 1} & \textbf{Métrica 2} \\
\midrule
\textit{Caso 1} & \textit{Valor} & \textit{Valor} \\
\textit{Caso 2} & \textit{Valor} & \textit{Valor} \\
\bottomrule
\end{tabular}
\end{table}

\subsection{Comparación}

\textbf{[Comparar la solución con una alternativa, una línea base o los
resultados esperados.]}

\subsection{Observaciones}

\textbf{[Registrar patrones, comportamientos inesperados o limitaciones
observadas durante las pruebas.]}

\section{Estructuras utilizadas}
\label{sec:Estructuras}
Para las implementaciones de la primera fase del proyecto se han utilizado circular queues y Circular linked lists.
\subsection{motivos de seleccion}
\begin{itemize}
\item Circular queue: Se decidio usar para el filtrado de señal al saber exactamente la cantidad de datos que se almacenaria y al tener un costo de tan solo O(1) para ingreso y salida de datos, ademas de que las FIFOS ecajaban adecuadamente con el proceso de filtrado
\item Circular linked list: Se uso para mantener un registro al ser datos de los cuales no se conoce la cantidad exacta de antemano, permitiendo un uso de memoria menos costoso
\end{itemize}

\subsection{Estructuras pensadas para futuras fases}
Por el momento la principal estructura pensada para una futura fase son los BST, estos se podria usar para por medio de umbrales de decision ser capaces de clasificar los latidos leidos del ECG.



\section{Primera implementacion}
\label{sec:prim_impl}
Como se ha mencionado a lo largo del documento se implemento un algoritmo basado en pan-Tompkins para el filtrado de señales ECG, su implementacion se realizo de la siguiente manera: 
\section*{Pseudocódigo del Algoritmo Pan-Tompkins}

\subsection{1. Etapas de Filtrado y Procesamiento}
En esta primera etapa se desarrollaron las funciones de filtrado de datos y procesamiento de muestras para despues ser aplicadas con la ADT
\begin{algorithm}[H]
\caption{Procesamiento de Muestra}
\begin{algorithmic}[1]
\Function{ProcesarMuestra}{$m$}
    \If{not tieneOffset}
        \State offset $\leftarrow$ $m$
        \State tieneOffset $\leftarrow$ true
    \EndIf
    \State $m \leftarrow m - offset$
    \State pb $\leftarrow$ \Call{PasaBajos}{$m$}
    \State pa $\leftarrow$ \Call{PasaAltos}{pb}
    \State derivar $\leftarrow$ \Call{Derivada}{pa}
    \State elevarCuadrado $\leftarrow$ \Call{Cuadrado}{derivar}
    \State ventana $\leftarrow$ \Call{IntegracionVentana}{elevarCuadrado}
    \State \Return ventana
\EndFunction
\end{algorithmic}
\end{algorithm}

\begin{algorithm}[H]
\caption{Filtro Pasa-Bajos}
\begin{algorithmic}[1]
\Function{PasaBajos}{$x$}
    \State Desplazar y almacenar $x$ en buffer Pasa-Bajos
    \State $y \leftarrow 2 \cdot y[n-1] - y[n-2] + x[n] - 2 \cdot x[n-6] + x[n-12]$
    \State Almacenar $y$ en buffer de salidas Pasa-Bajos
    \State \Return $y$
\EndFunction
\end{algorithmic}
\end{algorithm}

\begin{algorithm}[H]
\caption{Filtro Pasa-Altos}
\begin{algorithmic}[1]
\Function{PasaAltos}{$x$}
    \State Desplazar y almacenar $x$ en buffer Pasa-Altos
    \State $y \leftarrow y[n-1] - \frac{x[n]}{32.0} + x[n-16] - x[n-15] + \frac{x[n-32]}{32.0}$
    \State Almacenar $y$ en buffer de salidas Pasa-Altos
    \State \Return $y$
\EndFunction
\end{algorithmic}
\end{algorithm}

\begin{algorithm}[H]
\caption{Derivada}
\begin{algorithmic}[1]
\Function{Derivada}{$x$}
    \State Desplazar y almacenar $x$ en buffer de Derivada
    \State $y \leftarrow 0.125 \cdot (2 \cdot x[n] + x[n-1] - x[n-3] - 2 \cdot x[n-4])$
    \State \Return $y$
\EndFunction
\end{algorithmic}
\end{algorithm}

\begin{algorithm}[H]
\caption{Cuadrado e Integración de Ventana}
\begin{algorithmic}[1]
\Function{Cuadrado}{$x$}
    \State \Return $x \cdot x$
\EndFunction

\Function{IntegracionVentana}{$x$}
    \State datoAntiguo $\leftarrow$ Sacar el dato más antiguo del buffer de integración
    \State suma\_int $\leftarrow$ suma\_int + $x$ - datoAntiguo
    \State Desplazar y almacenar $x$ en buffer de integración
    \State \Return suma\_int / window\_size
\EndFunction
\end{algorithmic}
\end{algorithm}

\subsection{2. Fase de Decisión y Lógica de Umbrales}
En esta parte del algoritmo a partir de los umbrales se decide el tipo de entrada y se clasifica, actualizando los umbrales a partir de la informacion obtenida por este proceso.
\begin{algorithm}[H]
\caption{Fase de Decisión}
\begin{algorithmic}[1]
\Procedure{Decidir}{$v$}
    \State pendMax $\leftarrow$ $\max(\text{pendMax}, v - prev1)$
    \If{not inicializado}
        \State maxInicializacion $\leftarrow$ $\max(\text{maxInicializacion}, v)$
        \State sumaInicializacion $\leftarrow$ sumaInicializacion + $v$
        \If{$n + 1 \geq \text{muestrasInicializacion}$}
            \State spki $\leftarrow$ $0.25 \cdot \text{maxInicializacion}$
            \State npki $\leftarrow$ $0.5 \cdot (\text{sumaInicializacion} / \text{muestrasInicializacion})$
            \State \Call{ActualizarUmbrales}{}
            \State pendMax $\leftarrow$ $0$
            \State inicializado $\leftarrow$ true
        \EndIf
    \Else
        \If{$prev1 > prev2$ and $prev1 \geq v$}
            \State \Call{Clasificar}{$prev1, n - 1$}
            \State pendMax $\leftarrow$ $0$
        \EndIf
        \State \Call{Searchback}{}
    \EndIf
    \State $prev2 \leftarrow prev1$
    \State $prev1 \leftarrow v$
    \State $n \leftarrow n + 1$
\EndProcedure

\Procedure{ActualizarUmbrales}{}
    \State thresholdI $\leftarrow$ npki + $0.25 \cdot (spki - npki)$
    \State thresholdII $\leftarrow$ $0.5 \cdot thresholdI$
\EndProcedure
\end{algorithmic}
\end{algorithm}

\subsection{3. Clasificación de Picos y Searchback}
Este ultimo bloque del algoritmo se encarga de clasificar el tipo de pico a base de su pentiente y nuevamente actualizar los umbrales a partir de la informacion obtenida
\begin{algorithm}[H]
\caption{Clasificación de Picos y Searchback}
\begin{algorithmic}[1]
\Procedure{Clasificar}{pico, idx}
    \State periodoRefactor $\leftarrow$ $\text{round}(0.200 \cdot fs)$
    \State ventanaT $\leftarrow$ $\text{round}(0.360 \cdot fs)$
    \State desde $\leftarrow$ idx - ultimoQRS

    \If{ultimoQRS $\geq 0$ and desde $<$ periodoRefactor}
        \State \textbf{return}
    \EndIf

    \If{pico $>$ thresholdI}
        \State ondaT $\leftarrow$ (ultimoQRS $\geq 0$) and (desde $<$ ventanaT) and (pendMax $< 0.5 \cdot pendQRS$)
        \If{ondaT}
            \State npki $\leftarrow$ $0.125 \cdot pico + 0.875 \cdot npki$
        \Else
            \State spki $\leftarrow$ $0.125 \cdot pico + 0.875 \cdot spki$
            \State \Call{RegistrarQRS}{idx, pendMax}
        \EndIf
    \Else
        \State npki $\leftarrow$ $0.125 \cdot pico + 0.875 \cdot npki$
        \If{pico $>$ thresholdII and pico $>$ candPico}
            \State candPico $\leftarrow$ pico
            \State candIdx $\leftarrow$ idx
            \State candPend $\leftarrow$ pendMax
        \EndIf
    \EndIf

    \State \Call{ActualizarUmbrales}{}
\EndProcedure

\Procedure{Searchback}{}
    \If{candIdx $\geq 0$ and haySuficientesIntervalosRR and $(n - ultimoQRS) > 1.66 \cdot rrReferencia()$}
        \State spki $\leftarrow$ $0.25 \cdot candPico + 0.75 \cdot spki$
        \State \Call{RegistrarQRS}{candIdx, candPend}
        \State \Call{ActualizarUmbrales}{}
    \EndIf
\EndProcedure
\end{algorithmic}
\end{algorithm}


% =========================================================
% AGRADECIMIENTOS
% =========================================================
%ESTO ENTRA PARA POSTERIORES ENTREGAS
%\section*{Agradecimientos}

%Los autores agradecen a la Universidad Nacional de Colombia, Sede Amazonia,
%y a las personas o instituciones que apoyaron el desarrollo de este proyecto.

% =========================================================
% REFERENCIAS
% Sin BibTeX ni Biber para reducir tiempo de compilación.
% =========================================================

\begin{thebibliography}{00}

%\bibitem{cormen}
%T. H. Cormen, C. E. Leiserson, R. L. Rivest y C. Stein,
%\textit{Introduction to Algorithms}, 4th ed.
%Cambridge, MA, USA: MIT Press, 2022.

\bibitem{pt}
School of Electrical and Computer Engineering  Cornell David A. Duffield College of Engineering. Accedido el 14 de octubre de 2026. [En línea]. Disponible: \url{https://people.ece.cornell.edu/land/courses/ece5030/labs/s2014/Pan_Tompkins.pdf}

\bibitem{MIT-BIH}
G. Moody and R. Mark, "MIT-BIH Arrhythmia Database," PhysioNet, version 1.0.0, feb. 2005. doi: 10.13026/C2F305

\end{thebibliography}

\end{document}
